%% P12 --- Kombinatoryka i zliczanie
%%
%% Kazda liczba w tym programie, ktorej czytelnik nie policzy w pamieci,
%% pochodzi z code/p12_combinatorics.py i wchodzi na strone przez \val{}.
%% Listing konsolowy to zatwierdzony transcript pisany przez ten skrypt.
%%
%% TEZA: dwa pytania -- czy kolejnosc ma znaczenie i czy cos moze sie
%% powtorzyc -- rozstrzygaja, ktora z czterech liczb jest ta wlasciwa, a
%% niemal kazdy blad w zliczaniu bierze sie z odpowiedzi na zle pytanie.
%%
%% CO TEN PROGRAM DOSTAJE, odczytane z plikow, a nie z pamieci:
%%   F10  jest warstwa elementarna pod tym programem i sam to mowi. Daje
%%        regule iloczynu z warunkiem niezaleznosci, 2^n podzbiorow, liczbe
%%        par n(n-1)/2 z krokiem "podziel przez dwa" oraz wlaczenia i
%%        wylaczenia dla dwoch zbiorow -- a "cztery reguly formalnie,
%%        szufladki i rachunek urodzinowy" oddaje tutaj z nazwy. TRZY jego
%%        zatwierdzone wartosci sa w skrypcie sprawdzane bramkami.
%%   F04  daje silnie jako pi z indeksem w ciele oraz iloczyn pusty, przez co
%%        0! = 1 nie potrzebuje osobnej reguly.
%%   F03  daje logarytmy, jedyny sposob, by utrzymac 32000^20 w glowie.
%%   P02  ma sekcje "Nigdy nie twórz prawdopodobienstwa", a sekcja 4 tutaj to
%%        to samo odkrycie w miejscu, gdzie nikt sie go nie spodziewa:
%%        podrecznikowy wzor na kolizje zwraca DOKLADNIE ZERO we float64.
%%
%% CZEGO TEN PROGRAM NIE RUSZA, sprawdzone w tools/programs.json:
%%   czym JEST prawdopodobienstwo i co robi warunkowanie          -> P23
%%   zmienne losowe, wartosc oczekiwana, wariancja                -> P24
%%   grafy i ich argumenty zliczajace                             -> P13
%%   kwantyfikatory, indukcja, pisanie dowodu                     -> P14
%%   notacja O. Ten program liczy; nigdy nie mowi "kwadratowy"    -> P03
%%
%% Blizniak: programs/en/P12-combinatorics.tex. Ramka w ramke, makro w makro,
%% liczba w liczbe; porownuje je tools/parity.py.

\program{Kombinatoryka i zliczanie}
\label{prog:P12}
\index{zliczanie!cztery reguły}
\index{kombinatoryka}

\begin{outcomes}
\outcome{1--8}{zadać dwa pytania rozstrzygające, która reguła zliczania jest
  właściwa, i powiedzieć, dlaczego pytanie bez obu odpowiedzi nie jest jeszcze
  pytaniem}
\outcome{9--16}{policzyć wariacje i kombinacje oraz powiedzieć, skąd bierze
  się silnia w mianowniku symbolu Newtona}
\outcome{17--26}{stosować włączenia i wyłączenia dla więcej niż dwóch zbiorów
  oraz powiedzieć, co gwarantuje zasada szufladkowa, czego nie da żadne
  prawdopodobieństwo}
\outcome{27--35}{dobrać rozmiar skrótu do liczby elementów, które ma rozróżnić,
  i powiedzieć, który z dwóch wzorów na to wolno policzyć}
\outcome{36--43}{podać liczbę opisującą przestrzeń przeszukiwania i dokładną
  atrybucję oraz powiedzieć, co ta liczba wyklucza}
\end{outcomes}

\begin{quiz}
\item Z $\val{p12.four.n}$ cech wybierasz $\val{p12.four.k}$ do ablacji. Na ile
  sposobów? \teachesat{1--2}
  \answerto{Pytanie nie jest pełne. Dopóki nie powie, czy kolejność ma
  znaczenie i czy cechę wolno wybrać dwa razy, ma trzy dające się obronić
  odpowiedzi, a cała sekcja 1 to właśnie ta para pytań.}
\item Ile wynosi $\binom{\val{p12.sym.n}}{\val{p12.sym.k}}$ w porównaniu z
  $\binom{\val{p12.sym.n}}{\val{p12.sym.rest}}$? \teachesat{13--14}
  \answerto{Tyle samo, oba to $\val{p12.sym.val}$. Wybór
  $\val{p12.sym.k}$ cech do zostawienia jest wyborem
  $\val{p12.sym.rest}$ cech do odrzucenia, więc to jedno działanie opisane z
  dwóch stron.}
\item Trzy zbiory ewaluacyjne nakładają się parami, a część przykładów jest we
  wszystkich trzech. Zapisz rozmiar ich sumy. \teachesat{21--22}
  \answerto{Dodaj trzy, odejmij trzy nakładki parami, dodaj z powrotem część
  wspólną wszystkich trzech. Ten ostatni składnik jest tam dlatego, że
  odejmowanie par usunęło ją o raz za dużo.}
\item Skrót ma $\val{p12.bits64}$ bitów, a dokumentów jest
  $\val{p12.hash.m}$. Jaka jest z grubsza szansa, że dwa się zderzą?
  \teachesat{32--33}
  \answerto{Około $\val{p12.hash.p64}$ procent, czyli
  $\val{p12.hash.ratio}$ razy więcej niż liczba, do której dochodzi większość
  ludzi. Znaczenie ma liczba \emph{par}, a nie liczba dokumentów.}
\item Dokładne wartości Shapleya dla $\val{p12.shap.big}$ cech. Ile koalicji
  trzeba policzyć? \teachesatone{40}
  \answerto{$\val{p12.shap.big.coal}$, czyli $2^{\val{p12.shap.big}}$, a nie
  $\val{p12.shap.big}!$. Liczy się tylko to, \emph{które} cechy były
  wcześniej, nigdy ich kolejność, a zebranie kolejności zamienia jedną
  wykładniczą na mniejszą.}
\end{quiz}

% ============================================================
\section{Dwa pytania przed każdym wzorem}
\label{sec:P12-two-questions}
\index{zliczanie!kolejność i powtórzenia}

\begin{fr}
Część~IV to zmiana tematu, a nie ciąg dalszy. Nic z Programów
od~\ref{prog:P04} do~\ref{prog:P11} nie jest tu potrzebne i nic tutaj nie
potrzebuje macierzy.

Ten program dostaje w spadku Program~\ref{prog:F10}, który policzył trzy
rzeczy \dash{} przebiegi w siatce, podzbiory zbioru, pary z kolekcji \dash{} a
potem wprost nazwał to, co zostawia.

Zacznij tam, gdzie on skończył. Masz $\val{p12.four.n}$ cech i zamierzasz
poddać ablacji $\val{p12.four.k}$ z nich. Na ile sposobów można to zrobić?
\dotline
\nextframe
\end{fr}

\begin{fr}
\begin{ansblock}
Pytanie nie jest skończone, a jeśli podałeś liczbę, była jedną z trzech.
\end{ansblock}

Przed każdym wzorem dwa pytania:

\begin{itemize}
\item \textbf{Czy kolejność ma znaczenie?} Ablacja cech $3$, $7$ i $1$ w tej
  kolejności i ablacja tych samych cech w innej mogą być tym samym
  eksperymentem albo nie. Zdanie tego nie powiedziało.
\item \textbf{Czy coś może się powtórzyć?} Cechy nie da się poddać ablacji dwa
  razy w jednym przebiegu. Token można wylosować ze słownika dwa razy jak
  najbardziej. Tego zdanie też nie powiedziało.
\end{itemize}

\textbf{Niemal każdy błąd w zliczaniu w praktyce inżynierskiej to odpowiedź na
złe z tych dwóch pytań}, a tak łatwo je pominąć dlatego, że w zwykłej mowie
\enquote{wybierz trzy z dziesięciu} brzmi jak pełne polecenie.

Weź je po kolei. Załóżmy, że kolejność ma znaczenie, a powtórzenia są
dozwolone \dash{} losujesz $\val{p12.four.k}$ razy z
$\val{p12.four.n}$ i zapisujesz wynik. Ile jest wyników?
\nextframe
\end{fr}

\begin{fr}
\ans{$\val{p12.four.ordrep}$}
Czyli $\val{p12.four.n}^{\val{p12.four.k}}$, i jest to reguła iloczynu
z~Programu~\ref{prog:F10} bez żadnego dodatku: trzy niezależne wybory po
$\val{p12.four.n}$ możliwości.

Ogólnie \textbf{$n^{k}$}, i to jest liczba, której chcesz zawsze wtedy, gdy
liczony obiekt to \emph{ciąg z ustalonego alfabetu} \dash{} hasło, prompt o
$k$ tokenach, ciąg bajtów.

Teraz zakaż powtórzeń. Kolejność wciąż ma znaczenie, ale zużytej cechy już
nie ma. Ile teraz?
\dotline
\nextframe
\end{fr}

\begin{fr}
\ans{$\val{p12.four.ordnorep}$}
Z $\val{p12.four.n} \times 9 \times 8$: dziesięć na pierwszą, dziewięć na
drugą, bo jedna odpadła, osiem na trzecią.

Ta liczba ma nazwę \dash{} to liczba \textbf{wariacji bez powtórzeń}
$\val{p12.four.k}$ elementów z $\val{p12.four.n}$ \dash{} i warto się przy
niej zatrzymać, bo wygląda, jakby łamała regułę, która ją wyprodukowała.

Program~\ref{prog:F10} mówił, że niezależne wybory się mnożą, i starannie
powiedział, co znaczy niezależne: każdy wybór pierwszego zostawia te same
możliwości dla drugiego. Tu tak nie jest. Wybierz najpierw cechę $3$, a cecha
$3$ nie jest już dostępna jako druga. Czy reguła iloczynu obowiązuje, czy nie?
\nextframe
\end{fr}

\begin{fr}
\begin{ansblock}
Obowiązuje, a warunek, którego naprawdę potrzebuje, jest słabszy niż ten
podany.
\end{ansblock}

Ustalona musi być \textbf{liczba możliwości na każdym kroku, a nie same
możliwości}. Którakolwiek cecha poszła pierwsza, zostaje dokładnie dziewięć,
więc każdy z dziesięciu pierwszych wyborów rozgałęzia się na te same
dziewięć. Drzewo jest jednorodne, choć jego gałęzie noszą różne etykiety, a
jednorodne drzewo to wszystko, czego reguła iloczynu kiedykolwiek
potrzebowała.

\begin{note}
To doprecyzowanie reguły, którą już masz, a nie nowa reguła. Sformułowania
z~Programu~\ref{prog:F10} używaj przy sprawdzaniu siatki, bo tam możliwości
naprawdę są za każdym razem te same i jedna zakazana kombinacja psuje wynik.
To sformułowanie pozwala mnożyć mimo kurczącej się puli, a oba warto trzymać
osobno: pierwsze mówi o tym, które możliwości istnieją, drugie tylko o tym,
ile ich jest.
\end{note}

Teraz porzuć kolejność. Ablacja $\{3, 7, 1\}$ to jeden eksperyment, jakkolwiek
go wypiszesz. Ile teraz i przez co podzieliłeś?
\dotline
\nextframe
\end{fr}

\begin{fr}
\ans{$\val{p12.four.unordnorep}$, dzieląc przez $\val{p12.kfact}$}
Bo każdy wybór $\val{p12.four.k}$ cech był liczony raz na każdą kolejność, w
jakiej dało się go zapisać, a $\val{p12.four.k}$ elementów zapisuje się na
$\val{p12.four.k}! = \val{p12.kfact}$ sposobów.

\textbf{To jest krok \enquote{podziel przez dwa} z~Programu~\ref{prog:F10} w
wersji ogólnej.} Tamten program liczył pary, mnożąc $n$ przez $n-1$ i dzieląc
przez dwa, i mówił, że dzielenie jest faktem o zbiorach: $\{a, b\}$ i
$\{b, a\}$ to jedna para. To dzielenie było dzieleniem przez $2!$. Przy trzech
naraz jest to dzielenie przez $3!$, a przy $k$ naraz przez $k!$ \dash{} ten
sam krok, a jedyny powód, dla którego nie wyglądał na taki, jest ten, że
dzielenie przez dwa rzadko wygląda na dzielenie przez silnię.
\end{fr}

\begin{fr}
Zostaje jedna pusta komórka tabeli: kolejność nie ma znaczenia, a powtórzenia
\emph{są} dozwolone. Liczenie worków zamiast ciągów \dash{} ile różnych
multizbiorów $\val{p12.four.k}$ tokenów da się wylosować ze słownika o
$\val{p12.four.n}$ pozycjach, jeśli nie zapisujesz kolejności.

Odpowiedź to $\binom{n + k - 1}{k}$, czyli tutaj $\val{p12.four.unordrep}$.

\begin{rigourbox}
Ten wzór jest podany, a nie wyprowadzony. Argument nazywa się \emph{gwiazdki i
przegródki} i jest w każdym pierwszym kursie kombinatoryki: ustaw w rzędzie
$k$ elementów i $n-1$ przegródek, a każde ustawienie dwóch rodzajów symboli to
jeden multizbiór.

W praktyce jest najrzadziej używana z czterech i dlatego stoi tu jako nazwa do
rozpoznania, a nie jako technika. Kiedy już się pojawia, zwykle w przebraniu:
liczenie sposobów podziału ustalonego budżetu na $n$ kubełków to ta komórka, i
tak samo liczenie kształtów wektora worka słów o ustalonej sumie.
\end{rigourbox}
\end{fr}

\begin{fr}
Cała sekcja sprowadza się więc do jednej tabeli.

\begin{center}
\begin{tabular}{lll}
\toprule
 & \textbf{z powtórzeniami} & \textbf{bez powtórzeń} \\
\midrule
\textbf{kolejność ma znaczenie} & $n^{k}$ \dash{} $\val{p12.four.ordrep}$
  & $\frac{n!}{(n-k)!}$ \dash{} $\val{p12.four.ordnorep}$ \\
\textbf{kolejność bez znaczenia} & $\binom{n+k-1}{k}$ \dash{} $\val{p12.four.unordrep}$
  & $\binom{n}{k}$ \dash{} $\val{p12.four.unordnorep}$ \\
\bottomrule
\end{tabular}
\end{center}

Cztery liczby opisujące te same dziesięć rzeczy branych po trzy, biegnące od
$\val{p12.four.unordnorep}$ do $\val{p12.four.ordrep}$ \dash{} współczynnik
ponad ośmiu, rozstrzygnięty w całości przez dwa słowa, których w pytaniu nie
było.

\mermaidfig{p12-two-questions}{Dwa pytania wskazujące, której z czterech liczb
domaga się zadanie.}{kolejność i powtórzenia}
\end{fr}

% ============================================================
\section{Wariacje i kombinacje}
\label{sec:P12-choices}
\index{kombinatoryka!symbol Newtona}

\begin{fr}
Weź prawą górną komórkę do granicy. Zamiast ustawiać $\val{p12.four.k}$ z
$\val{p12.four.n}$ cech, ustaw je wszystkie.

Narzędzie już masz: Program~\ref{prog:F04} zapisał silnię jako pi z indeksem
we własnym ciele, $n! = \prod_{i=1}^{n} i$. Ile jest ustawień
$\val{p12.four.n}$ cech?
\nextframe
\end{fr}

\begin{fr}
\ans{$\val{p12.fact.n}$}
Czyli $\val{p12.four.n}!$, i to już liczba, której nikt nie wypisze.

Zauważ, jak wiąże się z liczbą wariacji. Ustawienie wszystkich dziesięciu to
ustawienie pierwszych trzech, a po nim ustawienie pozostałych siedmiu, więc
\[ n! = \frac{n!}{(n-k)!} \times (n-k)! \]
co jest kurczącą się pulą czytaną od tyłu i co tłumaczy, dlaczego liczbę
wariacji zapisuje się zwykle jako ten iloraz, a nie jako iloczyn $k$
malejących czynników. Iloraz jest formą schludniejszą; iloczyn tą, którą
policzysz w pamięci.
\end{fr}

\begin{fr}
Prawa dolna komórka ma symbol i nazwę, i warto mieć oba, bo tego używa
literatura.

\[ \binom{n}{k} = \frac{n!}{k!\,(n-k)!} \]

\textbf{symbol Newtona}: wariacje $k$ z $n$ podzielone przez $k!$ kolejności,
w których każdy wybór został policzony.

\begin{notationbox}
Po polsku czyta się to \emph{$n$ po $k$}, a sam symbol nazywa się symbolem
Newtona. Po angielsku ten sam obiekt czyta się \enquote{$n$ choose $k$}
\dash{} inna nazwa tego samego obiektu, a nie inny obiekt, i wydanie
angielskie tej książki używa nazwy angielskiej.

Sama notacja jest w obu ta sama, co jest ogólną zasadą tej książki: to, co
byś wpisał, zostaje bez zmian, to, co byś powiedział, niekoniecznie.
\end{notationbox}

Teraz użyj go do czegoś, co już policzyłeś inną drogą.
Program~\ref{prog:F10} policzył pary z $\val{p12.ie.a}$ dokumentów. Zapisz tę
liczbę jako symbol Newtona.
\dotline
\nextframe
\end{fr}

\begin{fr}
\ans{$\binom{\val{p12.ie.a}}{2} = \val{f10.pairs}$}
Para to dwa elementy wybrane z wielu z porzuceniem kolejności, czyli prawa
dolna komórka przy $k = 2$. Rozwiń go, a
$\frac{n!}{2!\,(n-2)!} = \frac{n(n-1)}{2}$, czyli wzór, który podał tamten
program.

\begin{note}
To nie jest podobieństwo. Skrypt \code{code/p12\_combinatorics.py} odczytuje
zatwierdzoną liczbę par z pliku wartości Programu~\ref{prog:F10} i przerywa
budowanie, jeśli $\binom{n}{2}$ jej nie odtworzy, więc oba programy nie mogą
rozjechać się co do tego, czym jest para. Dwie kolejne liczby tamtego programu
mają w sekcjach poniżej takie same bramki.
\end{note}
\end{fr}

\begin{fr}
Oto własność symbolu Newtona, w którą najtrudniej uwierzyć, i warto najpierw
się co do niej pomylić.

Z $\val{p12.sym.n}$ cech łatwiej jest wybrać $\val{p12.sym.k}$ z nich czy
wybrać $\val{p12.sym.rest}$ z nich? Odpowiedz szybko, bez liczenia.
\nextframe
\end{fr}

\begin{fr}
\begin{ansblock}
Ani jedno, ani drugie. Oba to $\val{p12.sym.val}$.
\end{ansblock}

\begin{trapbox}
\textbf{$\binom{n}{k} = \binom{n}{n-k}$, a powód mieści się w jednym zdaniu:
wybór $\val{p12.sym.k}$ cech do zostawienia \emph{jest} wyborem
$\val{p12.sym.rest}$ cech do odrzucenia.} Jest jedno działanie, a dwie liczby
opisują je z przeciwnych stron.

Powodem, dla którego niemal każdy wskazuje mniejszą liczbę, jest to, że oba
pytania sprawiają wrażenie różnej ilości pracy \dash{} wymienienie trzech
rzeczy idzie szybciej niż wymienienie siedemnastu. To fakt o zapisywaniu ich,
a nie o tym, ile ich jest. Wzór czyni to oczywistym, bo zamiana $k$ i $n-k$
zamienia dwa czynniki w mianowniku, który jest iloczynem.

Ciekawostką przestaje to być na końcach. $\binom{n}{0}$ i $\binom{n}{n}$ to
oba $1$ \dash{} jest dokładnie jeden sposób, by nie wybrać nic, i dokładnie
jeden, by wybrać wszystko \dash{} a pierwszy z nich to iloczyn pusty, za
którym Program~\ref{prog:F04} się wstawiał, przybywający tu jako liczba.
\end{trapbox}
\end{fr}

\begin{fr}
Liczby rosną więc od końców i mają szczyt pośrodku:
$\binom{\val{p12.sym.n}}{10} = \val{p12.sym.peak}$, największy wyraz w tym
wierszu.

Ten kształt rozstrzyga prawdziwy eksperyment. Masz $\val{p12.abl.n}$ cech i
chcesz zrobić badanie ablacyjne. Program~\ref{prog:F10} wycenia sprawdzenie
\emph{każdego} podzbioru na $2^{n}$, czyli $\val{p12.abl.all}$ przebiegów, i
to się nie wydarzy.

Ogranicz je zatem: żadnego podzbioru większego niż $\val{p12.abl.k}$. Ile to
przebiegów?
\dotline
\nextframe
\end{fr}

\begin{fr}
\ans{$\val{p12.abl.upto}$}
Z $1 + \val{p12.abl.singles} + \val{p12.abl.pairs} + \val{p12.sym.val}$
\dash{} pusta ablacja, pojedyncze, pary i trójki, czyli
$\binom{20}{0}$ do $\binom{20}{3}$, wśród nich $\val{p12.sym.val}$ sprzed
dwóch ramek.

\textbf{$\val{p12.abl.upto}$ przebiegów wobec $\val{p12.abl.all}$ to
współczynnik $\val{p12.abl.frac}$}, i to jest cała treść decyzji o
ograniczeniu rozmiaru podzbioru. To także powód, dla którego ogranicza się
zwykle do dwóch albo trzech, a nie do pięciu: wiersz wciąż tam stromo rośnie,
więc każdy kolejny rozmiar kosztuje więcej niż wszystkie mniejsze razem.

Nawyk jest ten z~Programu~\ref{prog:F10}, o regułę dalej: zapisz liczbę, zanim
napiszesz pętlę, a teraz umiesz ją zapisać dla każdego z czterech kształtów, a
nie tylko dla podzbiorów.
\end{fr}

% ============================================================
\section{Reguły, które nie są iloczynami}
\label{sec:P12-not-products}
\index{zliczanie!włączenia i wyłączenia}
\index{zliczanie!zasada szufladkowa}

\begin{fr}
Wszystko dotąd było regułą iloczynu w różnych przebraniach. Trzy rzeczy nią
nie są i każda odpowiada na pytanie, do którego iloczyn nie sięga.

Zacznij od najtańszej. Ustal jedną konkretną cechę \dash{} nazwij ją pierwszą
\dash{} i podziel każdy $k$-podzbiór zbioru $n$-elementowego na dwie kupki
według tego, czy ta cecha w nim jest. Policz obie kupki.
\nextframe
\end{fr}

\begin{fr}
\begin{ansblock}
$\binom{n-1}{k-1}$ ją zawiera, a $\binom{n-1}{k}$ nie, więc
\[ \binom{n}{k} = \binom{n-1}{k-1} + \binom{n-1}{k} \]
\end{ansblock}

Jeśli cecha jest w środku, pozostałe $k-1$ pochodzi z pozostałych $n-1$; jeśli
jej nie ma, wszystkie $k$ stamtąd pochodzą. Nic nie jest w obu kupkach i nic
nie jest poza obiema, więc liczby się dodają.

\textbf{To jest rekurencja i tak właśnie symbol Newtona się liczy.}
\code{math.comb} nie tworzy $n!$, by je potem wydzielić, co budowałoby liczbę
o setkach cyfr w odpowiedzi na pytanie, którego odpowiedź ma cztery; schodzi w
dół trójkąta dodawań, a każda wartość pośrednia jest liczbą czegoś.

Własny nawyk tej książki stosuje się i do wyprowadzenia, nie tylko do
arytmetyki: argument powyżej to \emph{decyzja tak-lub-nie o jednym elemencie},
czyli dokładnie ruch, którym Program~\ref{prog:F10} dostał $2^{n}$. Ten sam
ruch odpowiada tu na drugie pytanie, i to jest kolejna ramka.
\end{fr}

\begin{fr}
Zsumuj cały wiersz. Ile wynosi
$\binom{\val{p12.abl.n}}{0} + \binom{\val{p12.abl.n}}{1} + \cdots +
\binom{\val{p12.abl.n}}{\val{p12.abl.n}}$ i dlaczego?
\dotline
\nextframe
\end{fr}

\begin{fr}
\ans{$\val{p12.abl.all}$}
Czyli $2^{\val{p12.abl.n}}$, i jest to liczba podzbiorów
z~Programu~\ref{prog:F10} co do cyfry.

Powodem jest to, że wiersz \emph{jest} podzbiorami, poukładanymi według
rozmiaru. Żeby wybrać podzbiór, wolno wskazać jego rozmiar, a potem wskazać
elementy \dash{} to lewa strona \dash{} albo przejść przez elementy,
decydując: w środku czy poza, co jest stroną prawą. Jedna kolekcja obiektów
policzona dwa razy, a tożsamość wypada bez żadnej algebry.

\begin{note}
To druga z trzech bramek na Programie~\ref{prog:F10}. Skrypt sumuje ten wiersz
i przerywa budowanie, jeśli suma nie jest liczbą podzbiorów, którą tamten
program zatwierdził, więc \enquote{kombinacje} tutaj i \enquote{podzbiory}
tam nie mogą się rozjechać.
\end{note}
\end{fr}

\begin{fr}
Druga reguła, która nie jest iloczynem, jest u ciebie w wersji dla dwóch
zbiorów.

Program~\ref{prog:F10} miał dwa zbiory ewaluacyjne o $\val{p12.ie.a}$ i
$\val{p12.ie.b}$ przykładach, dzielące $\val{p12.ie.ab}$, i zwracał uwagę, że
dodanie rozmiarów liczy część wspólną dwa razy.

To są te same dwa zbiory. Przybył trzeci, o $\val{p12.ie.c}$ przykładach:
dzieli $\val{p12.ie.ac}$ z pierwszym i $\val{p12.ie.bc}$ z drugim, a
$\val{p12.ie.abc}$ przykładów jest we wszystkich trzech. Ile to różnych
przykładów?
\dotline
\nextframe
\end{fr}

\begin{fr}
\ans{$\val{p12.ie.union}$}
Nie $\val{p12.ie.naive}$, czyli o $\val{p12.ie.over}$ za dużo.

\[ \lvert A \cup B \cup C \rvert = \val{p12.ie.naive}
   - \val{p12.ie.pairsum} + \val{p12.ie.abc} = \val{p12.ie.union} \]

Dwa pierwsze składniki to reguła, którą miałeś. Trzeci jest nowy i to jego
warto zrozumieć, a nie zapamiętać: przykład należący do wszystkich trzech
zbiorów jest przez pierwszy składnik policzony \emph{trzy} razy, a potem przez
drugi odjęty \emph{trzy} razy, po jednym w każdej nakładce parami. Został więc
usunięty w całości i musi wrócić raz.

\textbf{Znaki się zmieniają na przemian, bo każda runda poprawek przestrzeliwa
poprzednią.} To reguła ogólna i właśnie dlatego nazywa się zasadą włączeń i
wyłączeń.

\begin{note}
Skrypt buduje te zbiory jako trzy prawdziwe zbiory z ich siedmiu obszarów i
porównuje wzór z rzeczywistą sumą, zamiast sprawdzać arytmetykę. Dwa pierwsze
zbiory ma też zabramkowane na zatwierdzonych rozmiarach
z~Programu~\ref{prog:F10}, więc jest to ten sam przykład ciągnięty dalej, a
nie nowy, który przypadkiem używa okrągłych liczb.
\end{note}
\end{fr}

\begin{fr}
Dla $n$ zbiorów reguła ma jeden składnik na każdy niepusty podzbiór zbiorów,
czyli $2^{n} - 1$ składników: $\val{p12.ie.terms3}$ dla trzech zbiorów i
$\val{p12.ie.terms10}$ dla $\val{p12.ie.sets}$.

\begin{trapbox}
\textbf{Nikt nie używa włączeń i wyłączeń powyżej trzech zbiorów, a powód
siedzi w tej liczbie.} Przy $\val{p12.ie.sets}$ źródłach danych żąda ona
$\val{p12.ie.terms10}$ przecięć, z których każde samo jest obliczeniem, i
każde musi być poprawne.

Ludzie robią zamiast tego sumę zbiorów i biorą jej rozmiar, co jest jednym
przejściem po danych i zerem składników poprawkowych \dash{} i jest to dobry
odruch, bo $2^{n}-1$ składników to $2^{n}$ podzbiorów bez pustego, więc reguła
kosztuje dokładnie tyle, ile wypisanie tego, czego wypisania miała oszczędzić.

Dwa i trzy zbiory to przypadki, które pojawiają się na przeglądzie projektu, i
oba mieszczą się w jednej linijce. Te zostaw; po resztę sięgnij po bibliotekę
zbiorów.
\end{trapbox}
\end{fr}

\begin{fr}
Trzecia reguła nie liczy niczego. Mówi, co \emph{musi} być prawdą, co jest
innym rodzajem zdania i okazuje się bardziej użyteczne.

Haszujesz $\val{p12.pig.items}$ dokumentów do $\val{p12.pig.buckets}$
kubełków. Powiedz coś, co jest na pewno prawdą o wyniku, jakakolwiek jest
funkcja skrótu i jakkolwiek dokumenty się rozłożą.
\dotline
\nextframe
\end{fr}

\begin{fr}
\ans{Któryś kubełek mieści co najmniej $\val{p12.pig.floor}$ dokumentów}
Bo gdyby każdy kubełek mieścił najwyżej $\val{p12.pig.floor} - 1$, suma
wynosiłaby najwyżej $\val{p12.pig.buckets} \times 3 = 768$, a dokumentów jest
$\val{p12.pig.items}$.

To \textbf{zasada szufladkowa}: przy $m$ elementach w $n$ szufladach któraś
mieści co najmniej $\lceil m/n \rceil$. Dowód jest ten podany wyżej i warto
zauważyć, że to dowód nie wprost, bez arytmetyki poza jednym mnożeniem.

\begin{note}
Skrypt sprawdza to na $\val{p12.pig.trials}$ losowych przydziałach i
\emph{sprawdza} jest tu właściwym słowem. Losowe przeszukanie nie może tu
niczego odkryć: kontrprzykład obaliłby dowód, więc to, co przeszukanie
kontroluje, to kod, dokładnie tak jak w~Programie~\ref{prog:P04}.
\end{note}

Daje też pierwszą uczciwą odpowiedź o kolizje skrótów. Przy większej liczbie
dokumentów niż kubełków kolizja nie jest prawdopodobna, jest
\textbf{pewna} \dash{} więc skrót o $\val{p12.bits64}$ bitach nie rozróżni
więcej niż $\val{p12.hash.n64}$ dokumentów, cokolwiek innego jest prawdą.
\end{fr}

\begin{fr}
A to ograniczenie jest niemal bezużyteczne, i o tym jest kolejna sekcja.

Nikt nie ma $\val{p12.hash.n64}$ dokumentów. Liczbą rozstrzygającą, czy skrót
jest dość duży, nie jest ta, przy której kolizje stają się pewne, tylko o
wiele mniejsza, przy której stają się \emph{prawdopodobne} \dash{} a o ile
mniejsza, to najbardziej zaskakująca liczba w tym programie.
\end{fr}

% ============================================================
\section{Ile bitów potrzebuje skrót}
\label{sec:P12-birthday}
\index{zliczanie!problem urodzinowy}
\index{haszowanie!kolizja}

\begin{fr}
Zacznij od wersji, którą każdy zna, bo pomylenie się w niej sprawia, że wersja
inżynierska trafia.

W pokoju jest $\val{p12.bday.m}$ osób. Jaka jest szansa, że dwie mają
urodziny tego samego dnia? Zgadnij, zanim policzysz; dni jest
$\val{p12.bday.n}$.
\dotline
\nextframe
\end{fr}

\begin{fr}
\ans{$\val{p12.bday.p}$ procent}
Raczej więcej niż połowa, przy $\val{p12.bday.m}$ osobach i
$\val{p12.bday.n}$ dniach.

Większość zgaduje coś w okolicach $\val{p12.bday.naive}$ procent, a to
zgadnięcie to $\val{p12.bday.m}$ podzielone przez $\val{p12.bday.n}$ \dash{}
czyli całkiem dobra odpowiedź na inne pytanie: na szansę, że ktoś ma urodziny
tego samego dnia co \emph{ty}.

\begin{trapbox}
\textbf{Pytanie nie brzmi: ty przeciw pokojowi. Brzmi: każda para w pokoju
przeciw każdej innej.} $\val{p12.bday.m}$ osób daje
$\binom{\val{p12.bday.m}}{2} = \val{p12.bday.pairs}$ par, a każda para to
własna szansa na trafienie.

$\val{p12.bday.pairs}$ szans po jednej na $\val{p12.bday.n}$ to jawnie niemała
liczba szans, i na tym polega cała niespodzianka. Liczba osób rośnie po
jednej; liczba par rośnie z kwadratem, co jest wzorem
z~Programu~\ref{prog:F10} i jego \enquote{dwa razy więcej danych to cztery
razy więcej pracy} przybywającym w zupełnie innym temacie.

Każdy, kogo to zaskakuje, jest zaskoczony z tego samego powodu: policzył
niewłaściwe obiekty. Policz pary, a nie zostaje nic do zaskoczenia.
\end{trapbox}

\mermaidfig{p12-pairs-not-items}{Dlaczego liczą się pary i co to robi z
potrzebnym rozmiarem skrótu.}{pary, nie elementy}
\end{fr}

\begin{fr}
Zdanie dokładne jest iloczynem. Nikt dotąd nie trafił, więc druga osoba musi
chybić pierwszej, trzecia musi chybić obu, i tak dalej:
\[ P(\text{no match}) = \prod_{i=0}^{m-1}\left(1 - \frac{i}{N}\right) \]
czyli pi z~Programu~\ref{prog:F04}, a zarazem rekurencja \dash{} każdy wyraz
to poprzedni razy jeden czynnik więcej.

Jest drugie wyrażenie na to samo i bierze się wprost z liczby par. Każda z
$\binom{m}{2}$ par chybia z prawdopodobieństwem $1 - 1/N$, więc
\[ P(\text{a match}) \approx 1 - e^{-m(m-1)/2N} \]
używając eksponenty z~Programu~\ref{prog:F07} do zamiany iloczynu liczb
bliskich jedynce na jeden wyraz.

Dwa wzory na jedną wielkość, jeden dokładny, drugi nie. Którego ma używać
kod?
\nextframe
\end{fr}

\begin{fr}
\begin{ansblock}
To zależy od reżimu i żaden nie jest bezpiecznym domyślnym wyborem.
\end{ansblock}

Zmierzone, jako błąd względny wobec odniesienia, którym nie jest żaden z nich:

\begin{center}
\begin{tabular}{llll}
\toprule
 & \textbf{iloczyn} & \textbf{eksponenta} \\
\midrule
$\val{p12.bday.m}$ osób, $\val{p12.bday.n}$ dni
  & $\val{p12.reg.bday.naive}$ & $\val{p12.reg.bday.closed}$ \\
$\val{p12.check.m}$ skrótów, $\val{p12.bits64}$ bitów
  & $\val{p12.reg.mid.naive}$ & $\val{p12.reg.mid.closed}$ \\
$\val{p12.trap.m}$ skrótów, $\val{p12.bits128}$ bitów
  & $\val{p12.reg.tiny.naive}$ & $\val{p12.reg.tiny.closed}$ \\
\bottomrule
\end{tabular}
\end{center}

\textbf{Każdy z nich jest dobry dokładnie tam, gdzie drugi zawodzi, i zawodzą
z przeciwnych powodów.} Iloczyn dokładny jest dokładny w przypadku
urodzinowym i zwraca błąd względny równy \emph{jeden} przy
$\val{p12.bits128}$ bitach, czyli nie zwraca nic. Eksponenta myli się o
$\val{p12.reg.bday.closed}$ w przypadku urodzinowym \dash{} daje
$\val{p12.bday.closed}$ procent tam, gdzie odpowiedź to
$\val{p12.bday.p}$ \dash{} a przy skali skrótów jest dobra do piętnastu cyfr.
\end{fr}

\begin{fr}
Błąd eksponenty łatwo wytłumaczyć: traktuje pary jako niezależne, a nie są
całkiem, i ma to znaczenie tylko wtedy, gdy jest ich dość mało, by to
zauważyć.

Awaria iloczynu to Program~\ref{prog:P02} przybywający tam, gdzie nikt się go
nie spodziewa. Oto ona w całości:

\transcript{p12-collision-zero}

Każdy czynnik $1 - i/N$ jest liczbą odległą od jedynki o mniej niż jedną
część na $10^{33}$, a Program~\ref{prog:P01} mówi, co float64 z takimi robi:
zaokrągla każdy dokładnie do $1$. Iloczyn jest więc dokładnie $1$, a odjęciu
go od $1$ nie zostaje nic do powiedzenia.

\begin{trapbox}
\textbf{Prawdopodobieństwo kolizji równe dokładnie zero czyta się jak dowód
bezpieczeństwa, a jest artefaktem zaokrąglenia.} Nie ma wyjątku, nie ma
ostrzeżenia i nie ma \code{nan} \dash{} funkcja zwraca liczbę
zmiennoprzecinkową, ta liczba to \code{0.0}, a \code{0.0} jest zupełnie
rozsądnie wyglądającą odpowiedzią na zadane pytanie.

Program~\ref{prog:P02} to nazwał i podał regułę: nigdy nie twórz
prawdopodobieństwa, odejmując od jedynki liczbę jej bliską. Wersja stąd jest
gorsza od tej, którą tamten program zmierzył, bo tam stratą były cyfry, a tu
są wszystkie, i dlatego, że zła odpowiedź jest tą, która kończy rozmowę.

Zwróć uwagę i na środkowy wiersz tabeli. Przy $\val{p12.bits64}$ bitach i
zaledwie $\val{p12.check.m}$ elementach iloczyn stracił \emph{już} pięć cyfr
znaczących, długo przed tym, nim zawiódł do końca. Wzór nie przestaje działać
przy progu; degraduje się, a degradacja jest niewidoczna, dopóki ktoś nie
policzy tego samego inną drogą.
\end{trapbox}
\end{fr}

\begin{fr}
Teraz pytanie inżynierskie. Deduplikujesz korpus przez haszowanie:
$\val{p12.hash.m}$ dokumentów, skrót o $\val{p12.bits64}$ bitach, czyli
$\val{p12.hash.n64}$ możliwych wartości.

Jaka jest szansa, że dwa różne dokumenty dostaną ten sam skrót i jeden
zostanie po cichu wyrzucony?
\dotline
\nextframe
\end{fr}

\begin{fr}
\ans{Około $\val{p12.hash.p64}$ procent}
A nie $\val{p12.hash.naive64}$, które sugeruje liczba dokumentów przez liczbę
wartości skrótu. \textbf{Te dwie różnią się o współczynnik
$\val{p12.hash.ratio}$}, a różnica to w całości pary.

\begin{aibox}
$\val{p12.hash.p64}$ procent to szansa, której nie zaakceptowałbyś nigdzie
indziej w potoku. To jeden korpus na czterdzieści zawierający co najmniej
jedną parę różnych dokumentów, którą krok deduplikacji uznał za identyczne
\dash{} a to, który z nich zostaje, jest przypadkowe, więc stratą jest
dokument, po cichu, bez linijki w logu i bez czegokolwiek dalej, co mogłoby to
zauważyć.

Powodem, dla którego $\val{p12.bits64}$ wydaje się z zapasem, jest rachunek,
którego nikt nie robi na głos: $\val{p12.hash.n64}$ to ogromna liczba, a
$\val{p12.hash.m}$ dokumentów to jej mały ułamek. Oba zdania są prawdziwe i
żadne nie jest tym pytaniem.
\end{aibox}
\end{fr}

\begin{fr}
Jak duży musi więc być korpus, żeby skrót o $\val{p12.bits64}$ bitach był
rzutem monetą?

Liczba par to około $m^{2}/2$, więc trafienie staje się prawdopodobne, gdy
$m^{2}/2N$ jest rzędu $1$, co stawia $m$ w okolicach $\sqrt{N}$. Zrobione
porządnie daje $m = \val{p12.half.c}\sqrt{N}$ dla dokładnie połowy.

\begin{center}
\begin{tabular}{lll}
\toprule
\textbf{skrót} & \textbf{wartości} & \textbf{elementów na rzut monetą} \\
\midrule
$\val{p12.bits64}$ bitów & $\val{p12.hash.n64}$ & $\val{p12.hash.half64}$ \\
$\val{p12.bits128}$ bitów & $\num{3.40e38}$ & $\val{p12.hash.half128}$ \\
\bottomrule
\end{tabular}
\end{center}

\textbf{Skrót o $\val{p12.bits64}$ bitach jest rzutem monetą przy
$\val{p12.hash.half64}$ dokumentach}, a taki korpus istnieje. Przy
$\val{p12.bits128}$ bitach ten sam próg to $\val{p12.hash.half128}$, a przy
$\val{p12.hash.m2}$ dokumentach prawdopodobieństwo kolizji wynosi
$\val{p12.hash.p128}$.

To jest cały argument za $\val{p12.bits128}$ bitami i nie jest nim to, że
$\val{p12.bits128}$ to dwa razy $\val{p12.bits64}$. \textbf{Podwojenie bitów
podnosi liczbę wartości do kwadratu, a więc podnosi do kwadratu korpus, który
zmieścisz}, bo bezpieczna pojemność idzie z pierwiastkiem: $2^{32}$ wobec
$2^{64}$.
\end{fr}

\begin{fr}
Reguła do wyniesienia z tej sekcji mieści się w jednej linijce i jest warta
więcej niż wzór, z którego pochodzi.

\textbf{Bezpieczna pojemność skrótu to mniej więcej pierwiastek z liczby
wartości, które może przyjąć}, bo różne muszą zostać pary, a tych jest około
$m^{2}/2$. Podziel wykładnik na pół i masz odpowiedź: $\val{p12.bits64}$
bitów mieści około $2^{32}$ elementów, $\val{p12.bits128}$ bitów około
$2^{64}$.

I dwa zdania, które ta sekcja powiedziała o kolizjach, warto trzymać osobno,
bo odpowiadają na różne pytania. Zasada szufladkowa mówi, że kolizje stają się
\textbf{pewne} przy $\val{p12.pig.certain64}$ elementach. Rachunek urodzinowy
mówi, że stają się \textbf{prawdopodobne} przy $\val{p12.hash.half64}$. Te
różnią się o współczynnik rzędu czterech miliardów i to zawsze ten drugi
rozstrzyga projekt.
\end{fr}

% ============================================================
\section{Dwie liczby, które rozstrzygają projekt}
\label{sec:P12-designs}
\index{zliczanie!przestrzeń przeszukiwania}
\index{zliczanie!wartość Shapleya}

\begin{fr}
Jeszcze dwie liczby i żadna nie potrzebuje niczego ponad to, co dały cztery
poprzednie sekcje. Potrzebują gotowości, by policzyć je przed zaproponowaniem
tego, co wyceniają.

Dekoder generuje $\val{p12.beam.len}$ tokenów ze słownika o
$\val{p12.beam.vocab}$ pozycjach. Ile różnych ciągów mógłby wyprodukować?
Przyda ci się Program~\ref{prog:F03}.
\dotline
\nextframe
\end{fr}

\begin{fr}
\ans{$\val{p12.beam.space}$}
Czyli $\val{p12.beam.vocab}^{\val{p12.beam.len}}$, lewa górna komórka tabeli
tego programu \dash{} kolejność ma znaczenie, a tokeny jak najbardziej się
powtarzają.

Wzięcie $\logb{10}$ z tego to jedyny sposób, by zobaczyć, co to jest:
$\val{p12.beam.len} \times \logb{10}\val{p12.beam.vocab}$, czyli ociupinę
ponad dziewięćdziesiąt. Liczba dziewięćdziesięciocyfrowa, przy dwudziestu
tokenach.

Teraz beam search o szerokości $\val{p12.beam.a}$ przechodzi po tej
przestrzeni. Na każdym z $\val{p12.beam.len}$ kroków przedłuża każdy ze swoich
$\val{p12.beam.a}$ żywych ciągów o każdy z $\val{p12.beam.vocab}$ tokenów i
zatrzymuje najlepsze $\val{p12.beam.a}$. Ile przedłużeń ocenia w sumie?
\nextframe
\end{fr}

\begin{fr}
\ans{$\val{p12.beam.scored.a}$}
Z $\val{p12.beam.len} \times \val{p12.beam.a} \times \val{p12.beam.vocab}$.
Jako ułamek przestrzeni jest to $\val{p12.beam.frac.a}$.

Podwój szerokość do $\val{p12.beam.b}$, a praca się podwaja, do
$\val{p12.beam.scored.b}$ ocenionych przedłużeń. Ułamek obejrzanej przestrzeni
też się podwaja, do $\val{p12.beam.frac.b}$.

\begin{warning}
\textbf{Ta liczba wyklucza argument, a nie praktykę.} Beam search w sposób
sprawdzalny daje lepsze wyjście niż dekodowanie zachłanne, a ta książka nie
zmierzyła, o ile ani gdzie przestaje pomagać.

Ta liczba rozstrzyga, \emph{dlaczego} nie może to być pokrycie. Przejście z
$\val{p12.beam.a}$ na $\val{p12.beam.b}$ przenosi ułamek obejrzanej
przestrzeni z $\val{p12.beam.frac.a}$ na $\val{p12.beam.frac.b}$, a oba są
zerem do każdego praktycznego celu. Szersza wiązka nie może pomagać przez
oglądanie większej części przestrzeni, bo żadna szerokość, którą ktokolwiek
uruchomi, nie ogląda jej wcale.

Cokolwiek daje, daje to z \emph{uporządkowania}, które model nakłada na
przestrzeń \dash{} dobre ciągi nie są rozrzucone po tych dziewięćdziesięciu
cyfrach losowo, tylko skupione tam, gdzie oceny samego modelu są wysokie, a
szersza wiązka zabezpiecza przed lokalnie mylącą oceną. To zdanie o modelu, a
nie o przeszukiwaniu, i to z nim należy dyskutować, gdy ktoś proponuje
podwojenie wiązki.
\end{warning}
\end{fr}

\begin{fr}
Ostatnia liczba to ta, która powinna zmienić sposób, w jaki czytasz wykres.

Wartość Shapleya przypisuje wyjście modelu jego cechom, a definicja jest
zdaniem zliczającym: \textbf{atrybucją cechy jest jej średni wkład krańcowy po
wszystkich kolejnościach cech.} Dodawaj cechy po jednej w jakiejś kolejności,
notuj, ile ta jedna dołożyła, gdy przyszła jej kolej, i uśrednij to po
wszystkich kolejnościach.

Przy $\val{p12.shap.big}$ cechach ile to kolejności?
\dotline
\nextframe
\end{fr}

\begin{fr}
\ans{$\val{p12.shap.big.ord}$}
Czyli $\val{p12.shap.big}!$, i to niczego nie rozstrzyga, bo tego nie da się
policzyć.

Ale spójrz jeszcze raz, czego definicja potrzebuje. Gdy przychodzi kolej
cechy, to, co dokłada, zależy od zbioru cech już obecnych \dash{} i zupełnie
nie zależy od kolejności, w jakiej je dodano. Każda kolejność stawiająca przed
nią ten sam zbiór daje więc ten sam wkład krańcowy i wolno je zebrać.

Zebranie zamienia $\val{p12.shap.big}!$ kolejności na
$2^{\val{p12.shap.big}} = \val{p12.shap.big.coal}$ podzbiorów, każdy z wagą
mówiącą, ilu kolejności był przedstawicielem. \textbf{Współczynnik
$\val{p12.shap.big.ratio}$}, za darmo, i tylko dzięki temu tę wartość w ogóle
da się policzyć.

\mermaidfig{p12-orderings-or-subsets}{Jak średnia po $n!$ kolejnościach zwija
się w sumę po $2^{n}$ podzbiorach.}{kolejności w podzbiory}
\end{fr}

\begin{fr}
\begin{note}
Tożsamość jest w \code{code/p12\_combinatorics.py} dowodzona przez wypisanie,
a nie stwierdzana: przy $\val{p12.shap.n}$ cechach skrypt liczy wartość każdej
cechy oboma sposobami \dash{} uśredniając po wszystkich
$\val{p12.shap.orderings}$ kolejnościach i sumując po wszystkich
$\val{p12.shap.coalitions}$ podzbiorach z wagami \dash{} i wymaga, by były
równe jako dokładne ułamki.

Ten sam sześciocechowy przykład pokazuje drugą rzecz, po którą sięga się do
atrybucji. Dwie z jego cech pokrywają dokładnie ten sam teren. Każda z nich
sama jest warta $\val{p12.shap.alone}$ i każdej przypisuje się
$\val{p12.shap.redundant}$, podczas gdy cesze pokrywającej teren, do którego
nic innego nie sięga, przypisuje się całe $\val{p12.shap.top}$, ile jest
warta. Atrybucje sumują się do $\val{p12.shap.total}$, które osiąga pełny
zbiór, i ta własność czyni z nich udziały, a nie oceny.
\end{note}

Zatem $\val{p12.shap.big.coal}$ obliczeń przy $\val{p12.shap.big}$ cechach, co
komputer zrobi. A teraz $\val{p12.shap.huge}$ cech. Ile?
\nextframe
\end{fr}

\begin{fr}
\ans{$\val{p12.shap.huge.coal}$}
Po milisekundzie na sztukę \dash{} co jest optymistyczne dla czegokolwiek z
modelem w środku \dash{} daje to $\val{p12.shap.huge.years}$ lat.

\begin{trapbox}
\textbf{Dokładna atrybucja Shapleya jest wykładnicza z konstrukcji, więc każda
implementacja, jakiej używałeś, jest przybliżeniem przez próbkowanie.} To nie
jest wada żadnej konkretnej biblioteki. Ta wielkość jest zdefiniowana jako
średnia po $2^{n}$ podzbiorach i nie ma przekształcenia, które zrobiłoby ich
mniej, bo każdy z tych podzbiorów naprawdę może dać inną odpowiedź.

To, co z tego wynika, warto zabrać na następny wykres atrybucji, jaki ci
pokażą. Słupki mają błąd próbkowania, tego błędu zwykle się nie rysuje i jest
on największy dokładnie tam, gdzie cechy oddziałują \dash{} czyli tam, gdzie
wykresu używa się do postawienia tezy. Dwa przebiegi tego samego wyjaśniacza
na tym samym wejściu potrafią przestawić środek rankingu.

Pytanie, które należy zadać takiemu wykresowi, jest tym, które ten program
zadaje od początku, i jest pytaniem o zliczanie: \textbf{ile z tych
$2^{n}$ faktycznie policzono?}
\end{trapbox}
\end{fr}

\begin{fr}
Jedna uwaga na koniec, i o to tak naprawdę w tym programie chodziło.

Każda liczba w nim była decyzją uczynioną widoczną. Czy skrót jest dość
szeroki, czy badanie ablacyjne zmieści się w tygodniu, czy da się poszerzyć
wiązkę, na ile można ufać atrybucji \dash{} każda z nich została
rozstrzygnięta liczbą, która zajmuje jedno mnożenie i której niemal nikt nie
liczy przed zobowiązaniem się do tego, co ona wycenia.

\textbf{Reguły nie są tu częścią użyteczną.} Jest ich cztery i mieszczą się na
karteczce. Użyteczny jest nawyk zauważania, że plan w ogóle zawiera jakąś
liczbę, a Program~\ref{prog:F10} ujął to najlepiej: zapisz liczbę, zanim
napiszesz pętlę.
\end{fr}

% ============================================================
\begin{summarybox}
\sumitem{1--2}{\textbf{Dwa pytania rozstrzygają, która reguła obowiązuje}: czy
  kolejność ma znaczenie i czy coś może się powtórzyć?
  \enquote{Wybierz $\val{p12.four.k}$ z $\val{p12.four.n}$} brzmi jak pełne
  polecenie i ma trzy dające się obronić odpowiedzi, dopóki obie nie zostaną
  ustalone.}
\sumitem{3--4}{Dwie z czterech liczb, dla $\val{p12.four.k}$ z
  $\val{p12.four.n}$: $n^{k} = \val{p12.four.ordrep}$ z kolejnością i
  powtórzeniami oraz $\frac{n!}{(n-k)!} = \val{p12.four.ordnorep}$ z
  kolejnością bez powtórzeń, gdzie pula kurczy się o jeden na każdym kroku.}
\sumitem{4--5}{Reguła iloczynu potrzebuje \textbf{ustalonej liczby możliwości
  na każdym kroku, a nie ustalonych możliwości}, i dlatego kurcząca się pula
  i tak się mnoży. Mocniejsze sformułowanie z~Programu~\ref{prog:F10} jest do
  sprawdzania siatki; to słabsze pozwala liczbie wariacji w ogóle być
  iloczynem.}
\sumitem{6}{$\binom{n}{k} = \val{p12.four.unordnorep}$ to wariacje podzielone
  przez $k! = \val{p12.kfact}$ kolejności, w których każdy wybór został
  policzony, a \textbf{to dzielenie jest krokiem \enquote{podziel przez dwa}
  z~Programu~\ref{prog:F10} w wersji ogólnej} \dash{} dzielił przez dwa, bo
  $2! = 2$.}
\sumitem{7--8}{Czwarta komórka, bez kolejności i z powtórzeniami, to
  $\binom{n+k-1}{k} = \val{p12.four.unordrep}$, podana i niewyprowadzona.
  Cztery liczby opisujące te same dziesięć rzeczy branych po trzy rozpinają
  współczynnik ponad ośmiu.}
\sumitem{9--10}{$\val{p12.four.n}!$ to $\val{p12.fact.n}$ ustawień, a
  $n! = \frac{n!}{(n-k)!} \times (n-k)!$ to kurcząca się pula czytana od tyłu
  \dash{} i dlatego liczbę wariacji zapisuje się jako iloraz, a liczy jako
  iloczyn.}
\sumitem{11--12}{$\binom{n}{k} = \frac{n!}{k!(n-k)!}$, po polsku \emph{symbol
  Newtona}. $\binom{n}{2} = \frac{n(n-1)}{2}$ to liczba par
  z~Programu~\ref{prog:F10}, a budowanie przerywa się, gdy te dwie się nie
  zgadzają.}
\sumitem{13--14}{\textbf{$\binom{n}{k} = \binom{n}{n-k}$, bo wybór tego, co
  zostaje, jest wyborem tego, co odpada.} Oba to $\val{p12.sym.val}$ przy
  $\val{p12.sym.n}$ branych po $\val{p12.sym.k}$ albo po
  $\val{p12.sym.rest}$.}
\sumitem{15--16}{Ablacja po każdym podzbiorze $\val{p12.abl.n}$ cech to
  $\val{p12.abl.all}$ przebiegów; ograniczona do $\val{p12.abl.k}$ cech to
  $\val{p12.abl.upto}$, czyli współczynnik $\val{p12.abl.frac}$.}
\sumitem{17--20}{\textbf{Rekurencja Pascala
  $\binom{n}{k} = \binom{n-1}{k-1} + \binom{n-1}{k}$} to rekurencja i tak
  właśnie symbol się liczy; suma wiersza daje
  $\val{p12.abl.all} = 2^{\val{p12.abl.n}}$, bo wskazanie rozmiaru, a potem
  elementów, to to samo działanie co decyzja w-środku-czy-poza dla każdego
  elementu.}
\sumitem{21--23}{\textbf{Włączenia i wyłączenia: znaki zmieniają się na
  przemian, bo każda runda poprawek przestrzeliwa poprzednią.} Trzy zbiory o
  $\val{p12.ie.a}$, $\val{p12.ie.b}$ i $\val{p12.ie.c}$ mieszczą
  $\val{p12.ie.union}$ różnych przykładów, a nie $\val{p12.ie.naive}$. Reguła
  ma $2^{n}-1$ składników, więc powyżej trzech zbiorów buduj sumę.}
\sumitem{24--26}{\textbf{Szufladki: $m$ elementów w $n$ szufladach wkłada co
  najmniej $\lceil m/n \rceil$ do którejś}, cokolwiek robi funkcja skrótu.
  Czyni kolizje \emph{pewnymi} przy $\val{p12.pig.certain64}$ elementach, a to
  nie jest liczba rozstrzygająca cokolwiek.}
\sumitem{27--28}{\textbf{Rachunek urodzinowy liczy pary, a nie elementy.}
  $\val{p12.bday.m}$ osób daje $\val{p12.bday.pairs}$ par i dzieli urodziny z
  prawdopodobieństwem $\val{p12.bday.p}$ procent, wobec
  $\val{p12.bday.naive}$ procent odpowiadających na inne pytanie.}
\sumitem{29--31}{Dwa wzory, każdy dobry tam, gdzie drugi zawodzi: iloczyn
  dokładny jest dokładny w przypadku urodzinowym i zwraca \textbf{dokładnie
  zero} przy $\val{p12.bits128}$ bitach, a eksponenta myli się tam o
  $\val{p12.reg.bday.closed}$ i jest dobra do piętnastu cyfr przy skali
  skrótów. Reguła Programu~\ref{prog:P02}, w funkcji skrótu.}
\sumitem{32--35}{\textbf{Bezpieczna pojemność skrótu to mniej więcej
  pierwiastek z liczby jego wartości.} $\val{p12.hash.m}$ dokumentów przy
  $\val{p12.bits64}$ bitach zderza się z prawdopodobieństwem
  $\val{p12.hash.p64}$ procent \dash{} $\val{p12.hash.ratio}$ razy więcej niż
  daje naiwne czytanie \dash{} a rzut monetą wypada przy
  $\val{p12.hash.half64}$ dokumentach wobec $\val{p12.hash.half128}$ dla
  $\val{p12.bits128}$ bitów.}
\sumitem{36--38}{$\val{p12.beam.vocab}^{\val{p12.beam.len}}$ to
  $\val{p12.beam.space}$ ciągów, a wiązka o $\val{p12.beam.a}$ ocenia
  $\val{p12.beam.scored.a}$ z nich. \textbf{Podwojenie wiązki podwaja ułamek
  $\val{p12.beam.frac.a}$}, więc pokrycie nie może być powodem, dla którego
  szersza wiązka pomaga.}
\sumitem{39--43}{\textbf{Wartość Shapleya uśrednia po każdej kolejności, a
  kolejności zwijają się w $2^{n}$ podzbiorów}, bo liczy się tylko to, które
  cechy były wcześniej. $\val{p12.shap.big.coal}$ przy
  $\val{p12.shap.big}$ cechach i $\val{p12.shap.huge.coal}$ przy
  $\val{p12.shap.huge}$, więc każda implementacja próbkuje \dash{} a słupki
  niosą błąd, którego nikt nie rysuje.}
\end{summarybox}

\canyou

\begin{testexercises}
\item Ze słownika o $50$ tokenach ile jest ciągów po $4$ tokeny? A ile
  \emph{zbiorów} $4$ różnych tokenów?
  \answerto{$50^{4} = \num{6250000}$ ciągów, bo kolejność ma znaczenie, a
  tokeny się powtarzają. I $\binom{50}{4} = \num{230300}$ zbiorów, dzieląc
  $50 \times 49 \times 48 \times 47$ wariacji przez $4! = 24$. Współczynnik
  ponad $27$ między nimi, rozstrzygnięty przez dwa słowa.}
\item Dwa źródła danych mieszczą $600$ i $900$ rekordów i dzielą $150$;
  trzecie mieści $400$, dzieląc $80$ z pierwszym, $120$ z drugim i $40$ z
  obydwoma. Ile to różnych rekordów?
  \answerto{$1900 - 350 + 40 = 1590$. Suma $600 + 900 + 400 = 1900$, odejmij
  $150 + 80 + 120 = 350$, dodaj z powrotem raz część wspólną wszystkich
  trzech, bo odjęcie trzech par usunęło ją trzy razy po tym, jak była
  policzona trzy razy.}
\item Haszujesz $10$ milionów dokumentów w przestrzeń $32$-bitową. Podaj
  przybliżone prawdopodobieństwo kolizji i powiedz, co czyni je tak dużym.
  \answerto{W praktyce $1$: liczba par to około $5 \times 10^{13}$ wobec
  $\num{4.29e9}$ wartości skrótu, więc wykładnik $m(m-1)/2N$ przekracza
  dziesięć tysięcy. Rzut monetą dla $32$ bitów wypada przy około
  $\num{7.7e4}$ dokumentach \dash{} $\val{p12.half.c}\sqrt{2^{32}}$ \dash{}
  więc dziesięć milionów jest daleko za nim.}
\item Siatka ma $6$ modeli, $4$ współczynniki uczenia i $3$ rozmiary wsadu,
  ale dwa modele nie uruchomią się przy największym wsadzie. Ile przebiegów?
  \answerto{$6 \times 4 \times 3 = 72$ minus zakazane. Dwa modele tracą po
  jednym rozmiarze wsadu przy wszystkich $4$ współczynnikach uczenia, więc
  odpada $8$: $64$ przebiegi. Warunek Programu~\ref{prog:F10} nałożony na
  regułę iloczynu i powód, by liczyć wykluczenia przy zapisywaniu siatki.}
\item Dlaczego $\binom{n}{k}$ liczy się rekurencją Pascala, a nie z silni?
  \answerto{Bo silnie są ogromne i skracają się niemal w całości:
  $\binom{50}{4}$ ma sześć cyfr, a $50!$ sześćdziesiąt pięć. Rekurencja
  Pascala schodzi w dół trójkąta dodawań, w którym każda wartość pośrednia
  sama jest liczbą czegoś, więc nic nie powstaje po to, by potem zostać
  wydzielone.}
\end{testexercises}

\begin{furtherproblems}
\item Pokaż, że $\binom{n}{k} = \binom{n}{n-k}$ ze wzoru, w jednej linijce.
  \answerto{$\binom{n}{n-k} = \frac{n!}{(n-k)!\,(n-(n-k))!} =
  \frac{n!}{(n-k)!\,k!}$, czyli $\binom{n}{k}$ z dwoma czynnikami mianownika
  zapisanymi w odwrotnej kolejności. Argument kombinatoryczny jest jeszcze
  krótszy: wybór $k$ do zostawienia jest wyborem $n-k$ do odrzucenia.}
\item Wiązka o szerokości $b$ przez $L$ kroków przy słowniku $V$ ocenia
  $LbV$ przedłużeń. Jaka szerokość byłaby potrzebna, by ocenić jeden procent
  przestrzeni, przy $V = \val{p12.beam.vocab}$ i $L = \val{p12.beam.len}$?
  \answerto{$b = \num{0.01}\,V^{L}/(LV)$, czyli około $2 \times 10^{84}$.
  Sensem zadania jest to, że ta liczba jest nie tyle wielka, co nieosiągalna,
  więc żadna szerokość nie jest argumentem z pokrycia.}
\item Iloczyn pusty daje $0! = 1$. Pokaż, że wymusza to rekurencja Pascala, a
  nie jest to umowa przyjęta dla wygody.
  \answerto{$\binom{n}{n} = \binom{n-1}{n-1} + \binom{n-1}{n}$, a ostatni
  składnik jest zerem, bo nie da się wybrać $n$ rzeczy z $n-1$. Zatem
  $\binom{n}{n} = \binom{n-1}{n-1} = \cdots = \binom{0}{0} = 1$, a
  $\binom{n}{n} = \frac{n!}{n!\,0!}$ równa się $1$ tylko wtedy, gdy $0! = 1$.
  Ta sama odpowiedź, do której Program~\ref{prog:F04} doszedł od akumulatora,
  od którego zaczyna pętla.}
\item Dlaczego podwojenie skrótu z $\val{p12.bits64}$ na
  $\val{p12.bits128}$ bitów podnosi mieszczony korpus do kwadratu, a nie
  podwaja go?
  \answerto{Bo bezpieczna pojemność idzie z $\sqrt{N}$, a podwojenie bitów
  podnosi $N$ do kwadratu: $\sqrt{N^{2}} = N$. Zatem $2^{32}$ staje się
  $2^{64}$, a wykładnik się podwaja, bo pierwiastek dzieli na pół podwojony
  wykładnik.}
\item Kolega proponuje policzyć dokładne wartości Shapleya dla modelu o
  $\val{p12.shap.huge}$ cechach, próbkując kolejności zamiast podzbiorów. Czy
  to pomaga?
  \answerto{Próbkowanie robi każda implementacja i jest jedynym wyjściem, ale
  czyni wynik oszacowaniem, a nie wartością dokładną \dash{} więc słowo
  \emph{dokładne} musi odpaść. Próbkowanie kolejności i próbkowanie
  podzbiorów to dwie drogi do tego samego estymatora, bo kolejności dzielące
  zbiór poprzedników cechy dają jeden wkład; wybór wpływa na wariancję, a nie
  na to, co się szacuje.}
\end{furtherproblems}
